\chapter{\texorpdfstring{$L^p$}{Lp} समष्टियाँ}\label{ch:b3:lp}

लेबेग समाकल विश्लेषण के लिए बनाया गया था; $L^p$ समष्टियाँ वे स्थान
हैं जहाँ वह विश्लेषण रहता है। वे बानाख समष्टियाँ हैं (रीस--फिशर)
--- अर्थात् वे \omterm{thm:b3:complete:completion}{पूर्णन} जिनका अभाव \cref{exo:b3:complete:1} ने \omterm{def:b3:topology:continuity}{संतत} फलनों में दिखाया था
--- और वे एक चिकनाई तकनीक का आधार देती हैं, \emph{मृदुकारकों के
साथ संवलन}, जो किसी भी $L^p$ फलन का $\mathcal C^\infty$ फलनों से आसन्नन कर देती है।
यह अध्याय होल्डर और मिन्कोव्स्की के समाकल रूप, \omterm{def:b3:complete:complete}{पूर्णता}, सघनता
प्रमेय, और नियमितीकरण मशीन सिद्ध करता है, और $L^p$ मापक्रम की
अंतर्वेशन एवं अंतर्गणन संबंधी भूगोल पर समाप्त होता है। सर्वत्र
$(X, \mathcal A, \mu)$ एक \omterm{def:b3:measure:measure}{माप समष्टि} है और फलन सम्मिश्र-मान वाले हैं; $\R^d$ पर \omterm{def:b3:measure:measure}{माप}
$\lambda_d$ है।

\section{परिभाषा; होल्डर और मिन्कोव्स्की}

\begin{definition}\label{def:b3:lp:space}
$1 \leq p < \infty$ के लिए $\mathcal L^p(\mu)$ $\norm f_p = \bigl(\int\abs
f^p\dd\mu\bigr)^{1/p} < \infty$ वाले \omterm{def:b3:lebesgue:measurable}{मापनीय} $f$ का समुच्चय है, और $\mathcal L^\infty(\mu)$
किसी शून्य समुच्चय के बाहर परिबद्ध $f$ का समुच्चय, जिस पर $\norm
f_\infty$
\emph{सारभूत उच्चतम} है --- अर्थात् वह न्यूनतम $M$ जिसके लिए
लगभग सर्वत्र $\abs f \leq M$ हो (यह निम्नतम प्राप्त हो जाता है: $M + \frac1n$ के लिए
शून्य समुच्चयों का प्रतिच्छेदन लीजिए)। चूँकि $\norm f_p = 0$ से केवल
\emph{लगभग सर्वत्र} $f
= 0$ अनिवार्य होता है (\cref{exo:b3:lebesgue:5}), हम
\[
L^p(\mu) = \mathcal L^p(\mu)/\{f = 0 \text{ लगभग सर्वत्र}\} :
\]
परिभाषित करते हैं; अवयव शून्य समुच्चयों के सापेक्ष फलनों के वर्ग
हैं, और $\norm\cdot_p$ $L^p$ पर एक सच्चा मानक है।\index{Lp समष्टि@$L^p$
समष्टि}
\end{definition}

\begin{theorem}[होल्डर की असमिका]\label{thm:b3:lp:holder}
मान लीजिए $\frac1p + \frac1q = 1$ (संयुग्मी घातांक) के साथ $1 \leq p, q \leq \infty$। \omterm{def:b3:lebesgue:measurable}{मापनीय} $f, g$ के
लिए:\index{होल्डर की असमिका}
\[
\norm{fg}_1 \leq \norm f_p\,\norm g_q ,
\]
और समता ($1 < p < \infty$, परिमित मानक, $f,g\ne0$ के लिए) तभी और केवल तभी होती है
जब $\abs f^p$ और $\abs g^q$ लगभग सर्वत्र समानुपाती हों।
\end{theorem}

\begin{proof}
$\{p, q\} = \{1, \infty\}$ की स्थितियाँ सीधी हैं (लगभग सर्वत्र $\abs{fg} \leq
\norm g_\infty\abs f$)। मान लीजिए $1 < p < \infty$;
$\norm f_p = \norm g_q = 1$ मानकीकृत कीजिए (समघातता; शून्य या अनंत मानक तुच्छ हैं)। यंग
की असमिका $ab \leq \frac{a^p}p +
\frac{b^q}q$ ($a, b \geq 0$; $\ln$ की अवतलता, जैसे \cref{pb:b3:banach:1} में) बिंदुशः
$\abs{f g} \leq
\frac{\abs f^p}p + \frac{\abs g^q}q$ देती है; समाकलन कीजिए: $\norm{fg}_1
\leq \frac1p + \frac1q = 1$। समता के लिए यंग में लगभग सर्वत्र
समता अनिवार्य है, अर्थात् लगभग सर्वत्र $\abs f^p = \abs g^q$ (मानकीकरण के बाद;
उसे उलटने पर समानुपातिता)।
\end{proof}

\begin{theorem}[मिन्कोव्स्की की असमिका]\label{thm:b3:lp:minkowski}
$1 \leq p \leq \infty$ के लिए: $\norm{f + g}_p \leq \norm f_p +
\norm g_p$।\index{मिन्कोव्स्की की असमिका}
\end{theorem}

\begin{proof}
$p = 1, \infty$: बिंदुशः/लगभग सर्वत्र त्रिभुज असमिका। $1 <
p < \infty$ के लिए मान लीजिए
$\norm{f+g}_p < \infty$ (अन्यथा $t^p$ की उत्तलता से मिलने वाली $\abs{f+g}^p \leq 2^{p-1}(\abs f^p + \abs g^p)$ का उपयोग करके
देखिए कि दायाँ पक्ष परिमित होने पर बायाँ भी परिमित है)। तब होल्डर
से
\[
\norm{f{+}g}_p^p \leq \int\abs f\,\abs{f{+}g}^{p-1} +
\int\abs g\,\abs{f{+}g}^{p-1}
\leq \bigl(\norm f_p + \norm g_p\bigr)\,
\bigl\|\abs{f{+}g}^{p-1}\bigr\|_q
\]
और $(p-1)q = p$ के कारण $\norm{\abs{f+g}^{p-1}}_q =
\norm{f+g}_p^{p/q}$; अब $\norm{f+g}_p^{p/q}$ से भाग दीजिए (यदि वह शून्येतर हो;
अन्यथा तुच्छ) और $p -
\frac pq = 1$ का उपयोग कीजिए।
\end{proof}

\section{पूर्णता और उसके सहचर}

\begin{theorem}[रीस--फिशर]\label{thm:b3:lp:complete}
$1 \leq p \leq \infty$ के लिए $L^p(\mu)$ एक बानाख समष्टि है। इसके अतिरिक्त $L^p$ में
अभिसरित होने वाले प्रत्येक अनुक्रम का कोई ऐसा उपानुक्रम है जो
\emph{लगभग सर्वत्र} अभिसरित होता है (और $p < \infty$ की स्थिति में किसी
$L^p$ प्रभावी फलन के साथ)।\index{रीस--फिशर प्रमेय}
\end{theorem}

\begin{proof}
$p = \infty$: कोई $\norm\cdot_\infty$-कोशी अनुक्रम एक ही शून्य समुच्चय (गणनीय संख्या के
शून्य समुच्चयों का सम्मिलन) के बाहर एकसमान कोशी होता है: अतः वह
उसके बाहर एकसमान रूप से अभिसरित होता है; काम पूरा। अब मान लीजिए
$p < \infty$। \cref{exo:b3:complete:1}(ख) से केवल निरपेक्षतः अभिसारी श्रेणियों का योग लेना
पर्याप्त है: मान लीजिए $\sum\norm{f_k}_p = M < \infty$। $G_n = \sum_{k \leq n}\abs{f_k}$ और $G = \sum_k\abs{f_k}$ रखिए ($[0,\infty]$ में
बिंदुशः): मिन्कोव्स्की से $\norm{G_n}_p \leq
M$, और एकदिष्ट अभिसरण प्रमेय ($G_n^p \nearrow G^p$)
$\int G^p \leq M^p$ दे देती है: अतः लगभग सर्वत्र $G < \infty$, इसलिए लगभग प्रत्येक $x$
के लिए श्रेणी $\sum f_k(x)$ निरपेक्षतः अभिसरित होती है; उसके योग को $S(x)$
कहिए (शून्य समुच्चय पर कोई भी मान)। तब $\abs{S - \sum_{k\leq n}f_k}^p \leq (2G)^p
\in L^1$, और प्रभावी अभिसरण
प्रमेय $\norm{S - \sum_{k\leq n}f_k}_p \to 0$ दे देती है: अर्थात् श्रेणी $L^p$ में अभिसरित होती है।

उपानुक्रम वाला कथन: यदि $L^p$ में $f_n \to f$ हो, तो $\norm{f_{n_{k+1}} - f_{n_k}}_p \leq 2^{-k}$ वाले $n_k$
चुनिए; श्रेणी $\sum(f_{n_{k+1}} - f_{n_k})$ पिछले तर्क के अंतर्गत आ जाती है: वह लगभग सर्वत्र
निरपेक्षतः अभिसरित होती है और किसी $G \in
L^p$ से प्रभावित है, अतः लगभग
सर्वत्र $f_{n_k} \to f_{n_1} + \sum(\cdots)$, और यह लगभग-सर्वत्र सीमा $f$ (का कोई प्रतिनिधि) ही
होनी चाहिए (दोनों $L^p$ सीमाएँ हैं)। प्रभावी फलन: $\abs{f_{n_k}} \leq \abs{f_{n_1}} + G$।
\end{proof}

\begin{remark}\label{rem:b3:lp:modes}
$L^p$ अभिसरण से लगभग सर्वत्र अभिसरण नहीं निकलता (\emph{टाइपराइटर}
अनुक्रम, \cref{exo:b3:lp:3}), और न ही विलोमतः (भागती उभारें): दोनों प्रकार केवल
उपानुक्रमों और प्रभुत्व के माध्यम से जुड़ते हैं। \cref{exo:b3:lp:3} के
प्रतिउदाहरण याद रखना सबसे अच्छा टीका है।
\end{remark}

\section{सघनता प्रमेय}

\begin{theorem}\label{thm:b3:lp:density}
मान लीजिए $1 \leq p < \infty$।
\begin{enumerate}
  \item \omterm{def:b3:lebesgue:simple}{सरल फलन} (परिमित \omterm{def:b3:measure:measure}{माप} वाले आलंब के साथ) $L^p(\mu)$ में सघन हैं।
  \item $L^p(\R^d)$ में \omterm{def:b3:topology:compact}{संहत} आलंब वाले \omterm{def:b3:topology:continuity}{संतत} फलन $\mathcal C_c(\R^d)$ सघन हैं।
  \item $L^p(\R^d)$ पर स्थानांतरण \omterm{def:b3:topology:continuity}{संतत} है: $\tau_hf = f(\cdot - h)$ लिखने पर $h \to 0$ जाने पर
        $\norm{\tau_hf - f}_p \to 0$।
\end{enumerate}
इनमें से कोई भी $p = \infty$ के लिए लागू नहीं होता।
\end{theorem}

\begin{proof}
(1) $f \geq 0$ के लिए: \cref{thm:b3:lebesgue:approximation} के द्विआधारी $s_n \nearrow f$ $\abs{f - s_n}^p
\leq f^p \in L^1$ को संतुष्ट करते
हैं: अतः प्रभावी अभिसरण प्रमेय। (प्रत्येक $s_n \leq f$ $L^p$ में है, और
जहाँ उसका मान धनात्मक है वहाँ उसके स्तर समुच्चयों का \omterm{def:b3:measure:measure}{माप} परिमित
है: $\mu(s_n \geq c) \leq c^{-p}\int f^p$।) व्यापक $f$ को चार अऋणात्मक भागों में बाँटिए।

(2) (1) से केवल $\lambda_d(A) < \infty$ वाले बोरेल $A$ के लिए $\mathbf 1_A$ का आसन्नन करना
पर्याप्त है। नियमितता (उपपत्ति \cref{thm:b3:measure:regularity} जैसी) $\lambda_d(U\setminus K) <
\varepsilon$ वाले \omterm{def:b3:topology:compact}{संहत} $K \subseteq A
\subseteq U$
और \omterm{def:b3:topology:topology}{विवृत समुच्चय} दे देती है; और उरिसोन का फलन
\[
\varphi(x) = \frac{d(x, \R^d\setminus U)}{d(x, \R^d\setminus U)
+ d(x, K)}
\]
\omterm{def:b3:topology:continuity}{संतत} है, $K$ पर $1$, $U$ के बाहर $0$, और उसे \omterm{def:b3:topology:compact}{संहत} आलंब
वाला लिया जा सकता है (पहले $U$ को किसी परिबद्ध \omterm{def:b3:topology:topology}{विवृत समुच्चय} तक
सिकोड़ दीजिए)। तब $\norm{\mathbf 1_A - \varphi}_p^p \leq \lambda_d(U\setminus K)
< \varepsilon$।

(3) $g \in \mathcal C_c$ के लिए: एकसमान \omterm{def:b3:topology:continuity}{सांतत्य} $\norm{\tau_hg - g}_\infty \to 0$ देता है, और $\abs h \leq 1$ के लिए
आलंब किसी स्थिर \omterm{def:b3:topology:compact}{संहत} समुच्चय में रहते हैं: अतः $\norm{\tau_hg - g}_p \to 0$। व्यापक $f$
के लिए: $\norm{f - g}_p <
\varepsilon$ वाला $g \in \mathcal C_c$ चुनिए; तब $\norm{\tau_hf - f}_p \leq 2\norm{f - g}_p +
\norm{\tau_hg - g}_p$ (मानक की स्थानांतरण
अपरिवर्तिता)।

$p = \infty$ के लिए: $\mathbf
1_{\intoo0\infty}$ का \omterm{def:b3:topology:continuity}{संतत} फलनों से एकसमान आसन्नन असंभव है
(उछाल), और $h \neq 0$ के लिए $\norm{\tau_h\mathbf 1_{\intoo0\infty} - \mathbf
1_{\intoo0\infty}}_\infty = 1$।
\end{proof}

\section{संवलन और नियमितीकरण}

\begin{theorem}[यंग की असमिका]\label{thm:b3:lp:young}
मान लीजिए $1 \leq p \leq \infty$, $f \in L^1(\R^d)$, $g \in
L^p(\R^d)$। तब $f * g$ लगभग सर्वत्र परिभाषित है,
$L^p$ में है, और\index{यंग की असमिका}
\[
\norm{f * g}_p \leq \norm f_1\,\norm g_p .
\]
\end{theorem}

\begin{proof}
$p = \infty$: सीधा परिबंध। $p = 1$: \cref{thm:b3:product:convolution}। अब मान लीजिए $1 < p < \infty$, $q$
संयुग्मी हैं। $\abs{f(y)} = \abs{f(y)}^{1/q}\cdot
\abs{f(y)}^{1/p}$ को बाँटिए और होल्डर लगाइए:
\[
\int\abs{f(y)}\,\abs{g(x{-}y)}\,\dd y
\leq \Bigl(\int\abs f\Bigr)^{1/q}
\Bigl(\int\abs{f(y)}\,\abs{g(x - y)}^p\,\dd y\Bigr)^{1/p} .
\]
$p$-वीं घात लीजिए और $x$ में समाकलन कीजिए; दूसरे गुणनखंड पर
टोनेली $\norm f_1^{p/q}\cdot\norm f_1\norm g_p^p$ दे देता है, अर्थात् $\norm{f*g}_p^p \leq \norm f_1^{1 + p/q}\norm g_p^p =
(\norm f_1\norm g_p)^p$ --- और टोनेली समाकल की
परिमितता \cref{thm:b3:product:convolution} की तरह लगभग सर्वत्र निरपेक्ष अभिसरण को उचित ठहरा
देती है।
\end{proof}

\begin{definition}[मृदुकारक]\label{def:b3:lp:mollifier}
फलन\index{मृदुकारक}\index{उभार फलन}
\[
\rho(x) = \begin{cases}
c\,\exp\Bigl(-\dfrac{1}{1 - \norm x^2}\Bigr) & \norm x < 1,\\
0 & \norm x \geq 1,
\end{cases}
\]
जिसमें $c$ $\int\rho = 1$ को मानकीकृत करता है, $\R^d$ पर $\mathcal C^\infty$ है: मर्म
यह है कि $t \mapsto \eu^{-1/t}\mathbf
1_{t>0}$ $\R$ पर $\mathcal C^\infty$ है और $0^+$ पर उसके सभी अवकलज $0$
हैं (प्रत्येक अवकलज किसी बहुपद $P$ के लिए $P(1/t)\eu^{-1/t}$ है, जो $0$ की
ओर जाता है; आगमन)। $\varepsilon > 0$ के लिए $\rho_\varepsilon(x) =
\varepsilon^{-d}\rho(x/\varepsilon)$ रखिए: यह $\bar B(0,
\varepsilon)$ में आलंबित है
और उसका समाकल अब भी $1$ है।
\end{definition}

\begin{theorem}[नियमितीकरण]\label{thm:b3:lp:regularization}
मान लीजिए $1 \leq p < \infty$ और $f \in L^p(\R^d)$। तब:
\begin{enumerate}
  \item $f * \rho_\varepsilon \in \mathcal C^\infty(\R^d)$, जिसमें $\partial^\alpha(f * \rho_\varepsilon) = f *
        \partial^\alpha\rho_\varepsilon$;
  \item $\varepsilon \to 0$ पर $\norm{f * \rho_\varepsilon - f}_p \to 0$;
  \item फलस्वरूप $\mathcal C^\infty_c(\R^d)$ $L^p(\R^d)$ में सघन है।
\end{enumerate}
\end{theorem}

\begin{proof}
(1) $x$ में समाकल के भीतर अवकलन (\cref{thm:b3:lebesgue:paramdiff}): किसी गोले $B$ में
$x$ के लिए $\abs{\partial_{x_i}\rho_\varepsilon(x - y)} \leq
C_\varepsilon\,\mathbf 1_{K}(y)$, जहाँ $K$ \omterm{def:b3:topology:compact}{संहत} है ($B$ से $\varepsilon$ के भीतर
$y$), और $\abs f\,\mathbf 1_K \in L^1$ ($\mathbf 1_K$ के सापेक्ष होल्डर): अतः प्रमेय लागू होती
है; ऊँचे अवकलजों के लिए दोहराइए।

(2) चूँकि $\int\rho_\varepsilon = 1$:
\[
(f * \rho_\varepsilon)(x) - f(x)
= \int \bigl(f(x - y) - f(x)\bigr)\rho_\varepsilon(y)\,\dd y ,
\]
और मिन्कोव्स्की की \emph{समाकल} असमिका से --- अथवा सीधे: प्रायिकता
\omterm{def:b3:measure:measure}{माप} $\rho_\varepsilon\dd y$ के साथ होल्डर/येंसन और टोनेली से ---
\[
\norm{f*\rho_\varepsilon - f}_p^p
\leq \int\Bigl(\int\abs{f(x-y) -
f(x)}^p\dd x\Bigr)\rho_\varepsilon(y)\,\dd y
= \int \norm{\tau_yf - f}_p^p\;\rho_\varepsilon(y)\,\dd y
\]
(बीच का चरण: भीतरी $y$-समाकल पर येंसन की असमिका \cref{exo:b3:lp:10} लगाइए,
फिर टोनेली)। समाकल्य $\norm y \leq \varepsilon$ में आलंबित है और $\varepsilon \to 0$ पर वहाँ एकसमान
रूप से $0$ की ओर जाता है (\cref{thm:b3:lp:density}(3)): अतः पूरा व्यंजक $0$ की ओर
जाता है।

(3) $f$ का $g \in \mathcal C_c$ से आसन्नन कीजिए (\cref{thm:b3:lp:density}(2)), फिर $g$ का $g *
\rho_\varepsilon \in \mathcal C_c^\infty$
से (\omterm{def:b3:topology:compact}{संहत} आलंब: आलंबों का योग)।
\end{proof}

\begin{example}[$\abs x$ का मृदुकरण, दरों सहित]
\label{ex:b3:lp:mollifyexample}
$\R$ पर $f(x) = \abs x$ लीजिए (जो स्थानीय रूप से $L^1$ है; प्रमेय प्रत्येक
परिबद्ध खिड़की पर लागू होती है) और कोई सममित \omterm{def:b3:lp:mollifier}{मृदुकारक} $\rho_\varepsilon$। तब
\[
f_\varepsilon(x) = (f * \rho_\varepsilon)(x)
= \int\abs{x - y}\,\rho_\varepsilon(y)\,\dd y
\]
$\mathcal C^\infty$ है; और मोड़ से दूर कुछ नहीं होता: $\abs x \geq \varepsilon$ के लिए $\rho_\varepsilon$ के आलंब
पर $\abs{x - y}$ $x$ में रैखिक है, अतः \emph{ठीक-ठीक} $f_\varepsilon(x) =
\abs x$ (सममिति
संशोधन को समाप्त कर देती है)। $0$ के निकट चिकनाई की कीमत ठीक
\[
0 \leq f_\varepsilon(0) = \int\abs
y\,\rho_\varepsilon(y)\,\dd y \leq \varepsilon,
\qquad
\norm{f_\varepsilon - f}_\infty \leq \varepsilon :
\]
है --- अर्थात् आसन्नन त्रुटि एकलता के $\varepsilon$-प्रतिवेश तक सीमित रहती
है और उसी के आकार की होती है। साथ ही सर्वत्र $f_\varepsilon'' \geq 0$ ($f$ उत्तल है,
और $\rho_\varepsilon
\geq 0$ के साथ संवलन उत्तलता को सुरक्षित रखता है), जिसमें $\int f_\varepsilon'' =
f_\varepsilon'(\infty) - f_\varepsilon'(-\infty) = 2$:
द्वितीय अवकलज द्रव्यमान $2$ की एक उभार है जिसे चौड़ाई $O(\varepsilon)$ में
दबा दिया गया है, अतः $\norm{f_\varepsilon''}_\infty \gtrsim
\varepsilon^{-1}$। चिकनाई एक सौदा है: एकसमान त्रुटि $O(\varepsilon)$
के बदले अवकलज का विस्फोट $O(\varepsilon^{-1})$ --- यही वह ठीक विनिमय दर है जिसे
मात्रात्मक विश्लेषण (अंतर्गणन असमिकाएँ, \cref{pb:b3:lp:1} के विचार-मंडल)
औपचारिक रूप देता है।
\end{example}

\begin{corollary}[विचरण कलन की मूल प्रमेयिका]\label{cor:b3:lp:fundlemma}
मान लीजिए $f \in L^1_{\mathrm{loc}}(\R^d)$ (संहतों पर \omterm{def:b3:lebesgue:l1}{समाकलनीय}) ऐसा है कि प्रत्येक $\varphi \in \mathcal
C^\infty_c(\R^d)$ के
लिए $\int f\varphi = 0$। तब लगभग सर्वत्र $f = 0$।\index{विचरण कलन की मूल
प्रमेयिका}
\end{corollary}

\begin{proof}
कोई गोला $B = B(0, R)$ स्थिर कीजिए और मान लीजिए $g = f\mathbf 1_{B(0, R+1)}
\in L^1$। $x \in B$ और $\varepsilon < 1$ के
लिए: $(g *
\rho_\varepsilon)(x) = \int f(y)\rho_\varepsilon(x - y)\dd y =
0$, जहाँ परीक्षण फलन $y \mapsto \rho_\varepsilon(x-y)
\in \mathcal C_c^\infty$ है। परंतु $L^1$ में $g * \rho_\varepsilon \to g$
(\cref{thm:b3:lp:regularization}): अतः $B$ पर लगभग सर्वत्र $g = 0$; अब $\R^d$ को निःशेष
कीजिए।
\end{proof}

\section{\texorpdfstring{$L^p$}{Lp} भूगोल}

\begin{proposition}\label{prop:b3:lp:inclusions}
(क) यदि $\mu(X) < \infty$ और $1 \leq p \leq q \leq \infty$ हों, तो $\norm f_p \leq
\mu(X)^{\frac1p - \frac1q}\,\norm f_q$ के साथ $L^q \subseteq L^p$।
(ख) $\R^d$ पर (अनंत \omterm{def:b3:measure:measure}{माप}) कोई अंतर्वेशन नहीं है: $p \neq q$ के लिए ऐसे
फलन हैं जो $L^p\setminus L^q$।
(ग) (अंतर्गणन) यदि $p < r < q$ हो और $\alpha \in \intoo01$ $\frac1r = \frac\alpha p + \frac{1 - \alpha}q$ से परिभाषित हो, तो
\[
\norm f_r \leq \norm f_p^{\alpha}\,\norm f_q^{1 - \alpha} ;
\]
और विशेष रूप से $L^p \cap L^q \subseteq L^r$।
\end{proposition}

\begin{proof}
(क) घातांकों $\frac qp$ और उसके संयुग्मी के साथ होल्डर: $\int\abs f^p\cdot 1 \leq \norm{\abs f^p}_{q/p}\,\norm
1_{(q/p)'} = \norm f_q^p\,\mu(X)^{1 - p/q}$ ($q = \infty$
सीधे)। (ख) $0$ के निकट और $\infty$ के निकट घातें $x^{-\alpha}$ अंशांकन कर
देती हैं: \cref{exo:b3:lp:2}। (ग) $\abs f^r = \abs
f^{r\alpha}\abs f^{r(1-\alpha)}$ लिखिए और संयुग्मी युग्म $\frac p{r\alpha}$, $\frac q{r(1-\alpha)}$
(जो $\alpha$ की परिभाषा से ठीक संयुग्मी हैं) के साथ होल्डर लगाइए:
$\int\abs f^r \leq \norm f_p^{r\alpha}\norm f_q^{r(1 -
\alpha)}$।
\end{proof}

\begin{method}\label{met:b3:lp:toolkit}
$L^p$ का औजार-बक्सा, जैसा आगे सर्वत्र प्रयुक्त होता है: समस्त
$f \in L^p$ के लिए कोई सर्वसमिका या असमिका सिद्ध करने के लिए --- उसे
किसी सघन वर्ग पर सिद्ध कीजिए (\cref{thm:b3:lp:regularization} के माध्यम से $\mathcal C_c^\infty$) और \omterm{def:b3:topology:continuity}{सांतत्य}
से बढ़ाइए (\cref{thm:b3:complete:extension}, क्योंकि दोनों पक्ष $L^p$-संतत हैं); $f = 0$ सिद्ध
करने के लिए $\mathcal
C_c^\infty$ के सापेक्ष परीक्षा लीजिए (\cref{cor:b3:lp:fundlemma}); चिकनाई पाने
के लिए संवलन कीजिए; और घातांक बदलने के लिए होल्डर तथा अंतर्गणन।
\cref{ch:b3:fouriertransform} का फूरिये सिद्धांत इसी विधि का एक लंबा अनुप्रयोग है।
\end{method}

\section{अभ्यास}

\begin{exercise}[$\star$]\label{exo:b3:lp:1}
(क) $L^2(\mu)$ में कोशी--श्वार्ज़ असमिका को होल्डर की $p = q = 2$ स्थिति के
रूप में कहिए और सिद्ध कीजिए।
(ख) किसी \emph{प्रायिकता} समष्टि पर दर्शाइए कि $p \mapsto \norm f_p$ अनवरोही है।
(ग) $p = 1$, $q = \infty$ के लिए होल्डर कब समता बनती है?
\end{exercise}

\begin{exercise}[$\star$]\label{exo:b3:lp:2}
किन $p \in \intco1\infty$ के लिए निम्नलिखित $L^p$ में हैं?
\[
x^{-1/2}\mathbf 1_{\intoo01},\qquad
x^{-1/2}\mathbf 1_{\intoo1\infty},\qquad
\frac{1}{x^{1/2}(1 + \abs{\ln x})}\ \intoo01 \text{ पर},
\qquad
\frac1{1 + \abs x}\ \R \text{ पर} .
\]
निष्कर्ष निकालिए: $\intoo01$ पर छोटा $p$ आसान है, $\intoo1\infty$ पर बड़ा
$p$ आसान है, और $\R$ पर कोई $L^p$ किसी दूसरे को अंतर्विष्ट नहीं
करता।
\end{exercise}

\begin{exercise}[$\star\star$]\label{exo:b3:lp:3}
(टाइपराइटर) द्विआधारी अंतरालों $I_1 =
\intcc01$, $I_2 = \intcc0{\frac12}$, $I_3 = \intcc{\frac12}1$, $I_4 = \intcc0{\frac14}$, \dots\
की गणना कीजिए और $f_n = \mathbf
1_{I_n}$ रखिए।
(क) दर्शाइए कि प्रत्येक $L^p(\intcc01)$, $p < \infty$ में $f_n \to 0$, पर \emph{प्रत्येक}
$x \in \intcc01$ के लिए $(f_n(x))$ अपसरित होता है।
(ख) \cref{thm:b3:lp:complete} द्वारा वचनित लगभग-सर्वत्र अभिसारी उपानुक्रम प्रस्तुत
कीजिए।
(ग) विलोमतः ऐसा अनुक्रम दीजिए जो लगभग सर्वत्र अभिसरित हो पर $L^1$
में नहीं, और ऐसा भी जो $L^1$ में अभिसरित हो पर किसी भी $L^p$,
$p > 1$ में नहीं।
\end{exercise}

\begin{exercise}[$\star\star$]\label{exo:b3:lp:4}
मान लीजिए $\mu(X) < \infty$ और $f \in L^\infty(\mu)$, $f \neq 0$। दर्शाइए कि $p \to \infty$ पर $\norm f_p \to \norm f_\infty$।
\emph{(ऊपरी परिबंध \cref{prop:b3:lp:inclusions} के (क) से; निचला परिबंध धनात्मक \omterm{def:b3:measure:measure}{माप} वाले
$\{\abs f > \norm f_\infty -
\varepsilon\}$ पर समाकलन करके।)}
\end{exercise}

\begin{exercise}[$\star\star$]\label{exo:b3:lp:5}
(क) \cref{thm:b3:lp:density}(3) की उपपत्ति $p < \infty$ का उपयोग ठीक कहाँ करती है?
(ख) दर्शाइए कि $f \in L^\infty(\R)$ $\norm{\tau_hf -
f}_\infty \to 0$ को संतुष्ट करता है तभी और केवल तभी जब
$f$ का कोई एकसमान \omterm{def:b3:topology:continuity}{संतत} प्रतिनिधि हो।
\end{exercise}

\begin{exercise}[$\star\star$]\label{exo:b3:lp:6}
मान लीजिए $p, q$ संयुग्मी हैं, $f \in L^p(\R^d)$, $g \in
L^q(\R^d)$। दर्शाइए कि $f * g$
\emph{सर्वत्र} परिभाषित है, परिबद्ध है, $\norm{f*g}_\infty \leq \norm f_p\norm g_q$ के साथ, और
\emph{एकसमान \omterm{def:b3:topology:continuity}{संतत}} भी। \emph{($L^p$ में स्थानांतरण का \omterm{def:b3:topology:continuity}{सांतत्य};
$p \in \{1, \infty\}$ की स्थिति अलग से लीजिए --- $p =
\infty$ के लिए $L^1$ गुणनखंड पर
स्थानांतरण \omterm{def:b3:topology:continuity}{सांतत्य} का उपयोग कीजिए।)}
\end{exercise}

\begin{exercise}[$\star\star$]\label{exo:b3:lp:7}
मान लीजिए $f \in L^1_{\mathrm{loc}}(\intoo ab)$ ऐसा है कि प्रत्येक $\varphi \in \mathcal
C^\infty_c(\intoo ab)$ के लिए $\int f\varphi' = 0$। दर्शाइए कि
$f$ लगभग सर्वत्र किसी अचर के बराबर है। \emph{($\int\chi = 1$ वाला $\chi \in \mathcal C_c^\infty$
स्थिर कीजिए; $\int\psi = 0$ वाला कोई भी $\psi \in \mathcal C_c^\infty$ एक $\varphi'$ है; किसी व्यापक
परीक्षण फलन को $\psi + (\int\psi)\chi$ के रूप में लिखिए और $c = \int
f\chi$ के साथ $f - c$ पर
\cref{cor:b3:lp:fundlemma} लगाइए।)}
\end{exercise}

\begin{exercise}[$\star\star\star$]\label{exo:b3:lp:8}
(चिकना उरिसोन) मान लीजिए $K \subseteq U \subseteq \R^d$, $K$ \omterm{def:b3:topology:compact}{संहत} है, $U$ विवृत।
$0 \leq \varphi \leq 1$, $K$ पर $\varphi = 1$, और $\operatorname{supp}\varphi \subseteq U$ वाला $\varphi \in \mathcal
C^\infty_c(\R^d)$ बनाइए। \emph{(छोटे
$\delta$ के लिए $K$ के $\delta$-प्रतिवेश $K_\delta$ के सूचक का $\rho_{\delta/2}$ से
मृदुकरण कीजिए।)} इससे परिमित संख्या के विवृत समुच्चयों से आच्छादित
किसी \omterm{def:b3:topology:compact}{संहत} समुच्चय के लिए $\mathcal C^\infty$ इकाई-विभाजन का कथन निकालिए।
\end{exercise}

\begin{exercise}[$\star\star$]\label{exo:b3:lp:9}
अंतर्गणन (\cref{prop:b3:lp:inclusions}(ग)) का उपयोग करते हुए:
(क) दर्शाइए कि सभी $p$ के लिए $\norm f_p \leq \norm f_1^{1/p}\norm
f_\infty^{1 - 1/p}$ के साथ $L^1(\R)\cap L^\infty(\R) \subseteq L^p(\R)$;
(ख) दर्शाइए कि स्थिर $f$ के लिए $f \mapsto \norm f_p$ $\frac1p$ में लघुगणकीय उत्तल
है, और ऐसा उदाहरण दीजिए जिसमें $f \in L^p$ ठीक किसी दिए हुए अंतराल
$(p_0, p_1)$ के $p$ के लिए हो।
\end{exercise}

\begin{exercise}[$\star\star$]\label{exo:b3:lp:10}
(येंसन) मान लीजिए $\mu$ एक \emph{प्रायिकता} \omterm{def:b3:measure:measure}{माप} है, $f \in
L^1(\mu)$
वास्तविक है, और $\Phi \colon \R \to \R$ उत्तल। दर्शाइए
\[
\Phi\Bigl(\int f\,\dd\mu\Bigr) \leq \int \Phi\circ
f\,\dd\mu
\]
\emph{(बिंदु $m = \int f$ पर $\Phi$ की आधार रेखा)}। इससे
समांतर--गुणोत्तर असमिका और \cref{exo:b3:lp:1}(ख) की $p \mapsto \norm f_p$ की एकदिष्टता निकालिए।
\index{येंसन की असमिका}
\end{exercise}

\begin{exercise}[$\star\star\star$]\label{exo:b3:lp:11}
(यंग की संवलन असमिका) मान लीजिए $\frac1p + \frac1q = 1 + \frac1r$ के साथ $1 \leq p, q, r \leq
\infty$, और $f \in
L^p(\R^d)$,
$g \in L^q(\R^d)$।
(क) $\norm{f * g}_r \leq \norm f_p\,\norm g_q$ सिद्ध कीजिए। \emph{($p, q,
r$ से निकाले गए संयुग्मी घातांकों
के लिए
\[
\abs{f(y)g(x-y)} = \bigl(\abs f^p\abs g^q\bigr)^{1/r}
\cdot\abs f^{\,p(1/p - 1/r)}\cdot\abs g^{\,q(1/q - 1/r)},
\]
लिखिए और घातांकों $r$, $\frac{pr}{r - p}$, $\frac{qr}{r - q}$ के साथ तीन-गुणनखंड वाली
होल्डर असमिका लगाइए; फिर टोनेली से $x$ में समाकलन कीजिए।)}
(ख) पहले से ज्ञात तीनों विशेष स्थितियाँ जाँचिए: $r =
\infty$ (होल्डर,
\cref{exo:b3:lp:6}); $q = 1$ (किसी \omterm{def:b3:lebesgue:l1}{समाकलनीय} नाभिक के साथ संवलन की $L^p$-स्थिरता);
और $p = q = 1$ ($L^1$ एक संवलन बीजगणित है, \cref{thm:b3:product:convolution})।
(ग) $\frac1p + \frac1q < 1 +
\frac1r$ के साथ कोई असमिका क्यों नहीं है? \emph{(प्रसार $f_\lambda(x) =
f(\lambda x)$ पर
परीक्षा लीजिए और दोनों पक्षों के मापन की तुलना कीजिए।)}
\end{exercise}

\begin{exercise}[$\star\star$]\label{exo:b3:lp:12}
(समता की स्थितियाँ) (क) होल्डर की असमिका $\int\abs{fg} \leq \norm f_p\norm g_q$ ($1 < p < \infty$) में
दर्शाइए कि समता तभी और केवल तभी होती है जब $\abs f^p$ और $\abs g^q$ लगभग
सर्वत्र समानुपाती हों। \emph{(यंग की असमिका $ab \leq \frac{a^p}p + \frac{b^q}q$ की समता स्थिति
का अनुसरण कीजिए, जो $a^p
= b^q$ है।)}
(ख) मिन्कोव्स्की की असमिका $\norm{f + g}_p \leq \norm f_p
+ \norm g_p$ ($1 < p < \infty$) में दर्शाइए कि $f,
g \neq 0$ के
साथ समता से लगभग सर्वत्र $c > 0$ के साथ $g = cf$ अनिवार्य हो जाता है।
(ग) $p = 1$ और $p = \infty$ से तुलना कीजिए: वहाँ की (बहुत बड़ी) समता
स्थितियों का उदाहरणों सहित वर्णन कीजिए।
\end{exercise}

\section{समस्या: हार्डी की असमिका}

\begin{problem}[{सप्ताहांत समस्या --- हार्डी की असमिका और उसका
तीक्ष्ण नियतांक}]\label{pb:b3:lp:1}
$f \in L^p(\intoo0{+\infty})$, $1 < p < \infty$ के लिए \emph{हार्डी संकारक}
\[
(Hf)(x) = \frac1x\int_0^x f(t)\,\dd t .
\]
परिभाषित कीजिए। हार्डी की असमिका (1920) कहती है\index{हार्डी की
असमिका}
\[
\norm{Hf}_p \;\leq\; \frac{p}{p-1}\,\norm f_p ,
\]
और नियतांक $\frac p{p-1}$ इष्टतम है तथा प्राप्त नहीं होता। यह समस्या सब
कुछ सिद्ध करती है, फिर श्रेणियों तक विस्तार करती है।

\medskip
\noindent\textbf{भाग I --- असमिका।} पहले मान लीजिए $f
\geq 0$ \omterm{def:b3:topology:continuity}{संतत} है
और $\intoo0{+\infty}$ में \omterm{def:b3:topology:compact}{संहत} आलंब वाला, और मान लीजिए $F(x) = \int_0^xf$।
\begin{enumerate}
  \item दर्शाइए कि $Hf \in L^p$: $0$ के निकट $F$ $0$ के किसी
        \omterm{def:b3:topology:topology}{प्रतिवेश} पर लुप्त हो जाता है; $\infty$ के निकट $F$ परिबद्ध
        है, अतः $(Hf)(x) = O(1/x)$, और $x \mapsto \frac1x$ $p > 1$ के लिए $L^p(\intoo1{+\infty})$ में है।
  \item खंडशः समाकलन करके दर्शाइए
        \[
        \int_0^\infty \Bigl(\frac Fx\Bigr)^p\dd x
        = \frac{p}{p-1}\int_0^\infty\Bigl(\frac
        Fx\Bigr)^{p-1}f(x)\,\dd x .
        \]
        \emph{($x^{1-p}F^p$ का अवकलन कीजिए; सीमांत पद लुप्त हो जाते हैं
        --- दोनों छोरों को उचित ठहराइए।)}
  \item दाईं ओर होल्डर लगाइए और ऐसे $f$ के लिए $\norm{Hf}_p \leq \frac p{p-1}\norm f_p$ निष्कर्ष
        निकालिए।
  \item समूचे $L^p$ तक बढ़ाइए: $f \geq 0$ के लिए $\intoo0{+\infty}$ में \omterm{def:b3:topology:compact}{संहत} आलंब
        वाला \omterm{def:b3:topology:continuity}{संतत} $f_n$ बनाइए जिसके लिए लगभग सर्वत्र $0 \leq f_n \nearrow f$
        \emph{(छाँटिए, फिर एकदिष्ट रूप से आसन्नन कीजिए --- रचना
        को उचित ठहराइए)}; तब बिंदुशः $Hf_n \nearrow Hf$ (औसत के भीतर एकदिष्ट
        अभिसरण प्रमेय) और एकदिष्ट अभिसरण प्रमेय असमिका को सीमा
        तक ले जाती है। चिह्नित या सम्मिश्र $f$ के लिए $\abs{Hf} \leq H\abs f$ से
        निष्कर्ष निकालिए।
\end{enumerate}

\medskip
\noindent\textbf{भाग II --- इष्टतमता।}
\begin{enumerate}[resume]
  \item $A > 1$ के लिए मान लीजिए $f_A(t) = t^{-1/p}\,\mathbf
        1_{\intcc1A}(t)$। $\norm{f_A}_p^p = \ln A$ परिकलित कीजिए, और
        $1 \leq x \leq A$ के लिए
        \[
        (Hf_A)(x) = \frac p{p-1}\;x^{-1/p}\,
        \bigl(1 - x^{-(1 - 1/p)}\bigr).
        \]
  \item $\liminf_{A\to\infty} \norm{Hf_A}_p/
        \norm{f_A}_p \geq \frac{p}{p-1}$ निष्कर्ष निकालिए, और निष्कर्ष निकालिए कि नियतांक
        इष्टतम है।
  \item दर्शाइए कि $f \neq 0$ के साथ समता $\norm{Hf}_p = \frac
        p{p-1}\norm f_p$ असंभव है। \emph{(प्रश्न
        3 में होल्डर की समता स्थिति का अनुसरण कीजिए: वह $f = cx^{-1/p}$
        प्रकार का व्यवहार अनिवार्य कर देती, जो $L^p$ में नहीं
        है।)}
\end{enumerate}

\medskip
\noindent\textbf{भाग III --- विविक्त असमिका।}
\begin{enumerate}[resume]
  \item $(0,\infty)$ पर किसी अनवर्धी $g \geq 0$ और $a_n = g(n)$ के लिए $\sum a_n^p$ और
        $\int g^p$ की तथा तदनुसार $H$-औसतों की तुलना करके भाग I से
        \emph{हार्डी की विविक्त असमिका} निकालिए: $a_n \geq 0$ के लिए
        \[
        \sum_{n\geq1}\Bigl(\frac{a_1 + \dots +
        a_n}{n}\Bigr)^{p}
        \;\leq\;
        \Bigl(\frac{p}{p-1}\Bigr)^{p}\,\sum_{n\geq1}a_n^p
        \]
        --- इसे पहले अनवर्धी $(a_n)$ के लिए उपर्युक्त तुलना से सिद्ध
        कीजिए, फिर व्यापक स्थिति को पुनर्विन्यास द्वारा अनवर्धी
        स्थिति तक सिमटाइए \emph{(एक पंक्ति के औचित्य के साथ यह
        स्वीकार कीजिए कि $(a_n)$ को घटते क्रम में छाँटने से बायाँ
        पक्ष केवल बढ़ सकता है जबकि दायाँ पक्ष अपरिवर्तित रहता
        है)}।
  \item निष्कर्ष निकालिए: यदि $\sum a_n^p < \infty$ हो, तो $(a_n)$ के चेज़ारो माध्य
        भी $\ell^p$ में होते हैं --- और ऐसा उदाहरण दीजिए ($p = 2$)
        जिसमें $(a_n) \in \ell^2$ हो पर $a_n$ योग्य न हो, फिर भी हार्डी माध्यों
        को नियंत्रित कर ले।
\end{enumerate}

\medskip
\noindent\textbf{भाग IV --- उपसंहार।}
\begin{enumerate}[resume]
  \item दर्शाइए कि $p = 1$ के लिए हार्डी की असमिका विफल हो जाती है:
        $f = \mathbf 1_{\intcc01}$ के साथ $Hf$ परिकलित कीजिए और $Hf \notin L^1$ पर ध्यान दीजिए।
        उपपत्ति कहाँ टूटती है?
\end{enumerate}

\medskip
\noindent\textbf{भाग V --- महत्तम फलन, और लेबेग की अवकलन
प्रमेय।} हार्डी मूल बिंदु से औसत लेता है; हार्डी--लिटिलवुड
\emph{प्रत्येक बिंदु के आसपास} औसत लेते हैं। $f \in L^1(\R)$ के लिए
परिभाषित कीजिए
\[
Mf(x) = \sup_{r>0}\ \frac1{2r}\int_{x-r}^{x+r}\abs
f\,\dd\lambda .
\]
\begin{enumerate}[resume]
  \item (विटाली, परिमित रूप) मान लीजिए $B_1, \dots, B_N$ विवृत अंतराल हैं।
        दर्शाइए कि कोई असंयुक्त उपकुल $B_{i_1}, \dots, B_{i_k}$ ऐसा है कि $\bigcup_jB_j
        \subseteq \bigcup_l3B_{i_l}$, जहाँ
        $3B$ उसी \omterm{ex:b3:groups:actions}{केंद्र} और तिगुनी लंबाई वाले अंतराल को निरूपित
        करता है \emph{(लोभी विधि: बार-बार वह सबसे लंबा अंतराल
        चुनिए जो पहले चुने गए अंतरालों से असंयुक्त हो)}।
  \item (दुर्बल प्रकार $(1,1)$) दर्शाइए कि प्रत्येक $t > 0$ के लिए
        \[
        \lambda\bigl(\{Mf > t\}\bigr) \;\leq\;
        \frac3t\,\norm f_1 :
        \]
        $Mf(x) > t$ वाला प्रत्येक $x$ $\int_{B_x}\abs f > t\,\lambda(B_x)$ वाला कोई केंद्रित अंतराल
        $B_x$ रखता है; कोई \omterm{def:b3:topology:compact}{संहत} $K \subseteq \{Mf > t\}$ लीजिए (आंतरिक नियमितता),
        उसे परिमित संख्या के $B_x$ से आच्छादित कीजिए, प्रश्न 11
        लगाइए, और निःशेष कीजिए।
  \item $x > 1$ के लिए $M\mathbf 1_{\intcc01}(x)$ परिकलित कीजिए और निष्कर्ष निकालिए कि
        प्रत्येक $f \neq 0$ के लिए $Mf \notin L^1$ (अनंत पर $Mf(x) \geq \frac c{\abs x}$): अर्थात्
        $p =
        1$ पर प्रश्न 12 की दुर्बल असमिका ही सर्वोत्तम संभव
        कथन है।
  \item ($p > 1$ के लिए प्रबल प्रकार) $f \in L^p$ के लिए: $f =
        f\,\mathbf 1_{\abs f > t/2} + f\,\mathbf 1_{\abs f
        \leq t/2}$ को
        बाँटिए, $Mf \leq M\bigl(f\mathbf 1_{\abs
        f > t/2}\bigr) + \frac t2$ पर ध्यान दीजिए, और प्रश्न 12 को परत-केक
        सूत्र (\cref{prop:b3:product:layercake}) तथा टोनेली के साथ मिलाकर सिद्ध कीजिए
        \[
        \norm{Mf}_p^p \;\leq\;
        \frac{6p\,2^{p-1}}{p-1}\,\norm f_p^p .
        \]
        ($p \downarrow 1$ पर विस्फोट प्रश्न 13 की विफलता ही है, मात्रात्मक
        रूप में।)
  \item (लेबेग की अवकलन प्रमेय) सिद्ध कीजिए: $f \in
        L^1(\R)$ के लिए
        \[
        \frac1{2r}\int_{x-r}^{x+r}f\,\dd\lambda
        \;\longrightarrow\; f(x)
        \qquad (r \to 0)\quad\text{लगभग प्रत्येक }x\text{ के लिए} .
        \]
        \emph{(\omterm{def:b3:topology:continuity}{संतत} $f$ के लिए स्पष्ट। व्यापक स्थिति में $f
        = g + h$
        लिखिए, जहाँ $g$ \omterm{def:b3:topology:compact}{संहत} आलंब वाला \omterm{def:b3:topology:continuity}{संतत} है और
        $\norm h_1 < \varepsilon$
        (\cref{thm:b3:lp:density}); जिस समुच्चय पर औसतित दोलन का $\limsup_{r\to0}$ $\delta$ से अधिक
        है वह $\{Mh > \delta/2\} \cup \{\abs h
        > \delta/2\}$ के भीतर बैठता है, जिसका \omterm{def:b3:measure:measure}{माप} $O(\varepsilon/\delta)$ है; अब
        $\varepsilon \to 0$ लीजिए, फिर किसी अनुक्रम के अनुदिश $\delta \to 0$।)}
  \item निष्कर्ष निकालिए: (क) लगभग प्रत्येक बिंदु $f$ का
        \emph{लेबेग बिंदु} है; (ख) $f \in L^1$ के लिए आद्यंतर $F(x) = \int_0^xf$
        लगभग सर्वत्र अवकलनीय है और लगभग सर्वत्र $F' =
        f$ --- अर्थात्
        लेबेग संसार में कलन की मूल प्रमेय का समाकल वाला आधा भाग,
        जो सीढ़ी (\cref{pb:b3:measure:1}) से खुले वृत्त को बंद कर देता है, क्योंकि
        सीढ़ी ने दिखाया था कि विलोम आधा भाग विफल हो सकता है।
  \item (\omterm{ex:b3:lebesgue:gamma}{घनत्व} बिंदु) \omterm{def:b3:lebesgue:measurable}{मापनीय} $A \subseteq \R$ के लिए दर्शाइए कि लगभग
        प्रत्येक $x \in A$ $\frac{\lambda(A\cap\intcc{x-r}{x+r})}{2r} \to 1$ को संतुष्ट करता है: अर्थात् \omterm{def:b3:lebesgue:measurable}{मापनीय}
        समुच्चय अपने लगभग सभी बिंदुओं पर स्थानीय रूप से भरे हुए
        होते हैं। दो पंक्तियों में रेखांकित कीजिए कि इससे
        श्टाइनहाउस की प्रमेय (\cref{exo:b3:measure:8}) की एक और उपपत्ति कैसे निकलती
        है।
\end{enumerate}

\medskip
\noindent\textbf{भाग VI --- औसतीकरण की विषय-वस्तु पर विविधताएँ।}
\begin{enumerate}[resume]
  \item (भारित हार्डी) $\alpha < p - 1$ और $f \geq 0$ के लिए उसी खंडशः समाकलन से
        दर्शाइए
        \[
        \int_0^\infty\Bigl(\frac{F(x)}x\Bigr)^{p}
        x^{\alpha}\,\dd x \;\leq\;
        \Bigl(\frac{p}{p - 1 - \alpha}\Bigr)^{p}
        \int_0^\infty f(x)^p\,x^{\alpha}\,\dd x
        \]
        और जाँचिए कि सीमांत स्थिति $\alpha = p - 1$ वास्तव में वर्जित है
        (प्रश्न 10 का प्रतिउदाहरण अनुकूलित कीजिए)।
  \item (संलग्न) मान लीजिए $H^*f(x) =
        \int_x^{\infty}\frac{f(t)}t\,\dd t$। अऋणात्मक $f, g$ के लिए $\langle
        Hf, g\rangle = \langle f, H^*g\rangle$
        दर्शाइए (टोनेली), और $\norm{H^*f}_p \leq p\,\norm f_p$ सिद्ध कीजिए \emph{(सीधे खंडशः,
        अथवा द्वैतता से संयुग्मी घातांक पर हार्डी से --- ध्यान
        दीजिए कि कौन-सा घातांक कौन-सा नियतांक उठाता है)}।
  \item (हिल्बर्ट प्रकार की एक असमिका) इससे निष्कर्ष निकालिए कि
        अऋणात्मक $f \in L^p$, $g \in L^q$ के लिए:
        \[
        \int_0^\infty\!\!\int_0^\infty
        \frac{f(x)\,g(y)}{\max(x,y)}\,\dd x\,\dd y
        \;\leq\; (p + q)\,\norm f_p\,\norm g_q
        \]
        \emph{($y < x$ / $y \geq x$ के अनुदिश बाँटिए: प्रत्येक आधा एक
        फलन की दूसरे के हार्डी रूपांतर के साथ युग्मन है)}।
  \item (इष्टतमता, विविक्त) दर्शाइए कि प्रश्न 8 का नियतांक $\bigl(\frac p{p-1}\bigr)^p$
        भी इष्टतम है: $a_n = n^{-1/p}\,\mathbf 1_{n \leq
        N}$ पर परीक्षा लीजिए, दोनों पक्षों की
        समाकलों से तुलना कीजिए, और $N
        \to \infty$ जाने दीजिए (भाग II का
        विविक्त दर्पण)।
  \item (संश्लेषण) इस समस्या में तीन औसतीकरण संकारक आए: हार्डी का
        $H$, विविक्त चेज़ारो माध्य, और महत्तम संकारक $M$।
        प्रत्येक की परिबद्धता क्या कहती है, यह एक-एक पंक्ति में
        कहिए, ध्यान दीजिए कि तीनों ठीक $p = 1$ पर विफल होते हैं, और
        समझाइए कि यह एक ही विफलता तीन बार क्यों है (हरात्मक पुच्छ
        $\frac1x$)।
\end{enumerate}

\medskip
\noindent\textbf{भाग VII --- कार्लमान की असमिका, और तीक्ष्ण
कितना तीक्ष्ण।}
\begin{enumerate}[resume]
  \item (कार्लमान की असमिका) मान लीजिए $\sum
        a_n < \infty$ के साथ $a_n \geq 0$। प्रश्न
        8 की विविक्त हार्डी असमिका $b_n = a_n^{1/p}$ पर लगाइए,
        समांतर--गुणोत्तर माध्य असमिका का उपयोग कीजिए, और $p \to
        \infty$
        जाने दीजिए (दर्शाइए कि $p \mapsto
        \bigl(\frac p{p-1}\bigr)^p$ घटकर $\eu$ तक जाता है),
        जिससे
        \[
        \sum_{n\geq1}\bigl(a_1a_2\cdots a_n\bigr)^{1/n}
        \;\leq\; \eu\,\sum_{n\geq1}a_n :
        \]
        मिल जाए: अर्थात् किसी योग्य अनुक्रम के गुणोत्तर माध्य भी
        योग्य होते हैं, और कीमत अधिकतम $\eu$।
  \item (नियतांक $\eu$ इष्टतम है) $a_n =
        \frac1n\,\mathbf 1_{n\leq N}$ पर परीक्षा लीजिए: \cref{pb:b3:product:1} के
        स्टर्लिंग कोष्ठन का उपयोग करके $(n!)^{-1/n} = \frac\eu n\bigl(1 +
        O\bigl(\frac{\ln n}n\bigr)\bigr)$ दर्शाइए, निष्कर्ष
        निकालिए कि कार्लमान के दोनों पक्ष $\eu\ln N$ की तरह बढ़ते हैं,
        और निष्कर्ष निकालिए कि $\eu$ से छोटा कोई नियतांक काम नहीं
        कर सकता। (प्रतिरूप पर ध्यान दीजिए: हार्डी और कार्लमान
        दोनों के इष्टतमकारी वही हरात्मक प्रकार के अनुक्रम हैं जो
        समष्टि में होने से बाल-बाल चूक जाते हैं।)
  \item (``तीक्ष्ण'' तक कितनी धीमी पहुँच?) $p = 2$ लीजिए। $f = \mathbf 1_{\intcc01}$ के
        लिए परिबंध $2$ के सामने $\norm{Hf}_2/\norm f_2 = \sqrt2$ परिकलित कीजिए। प्रश्न 5
        के निकट-इष्टतमकारी $f_A$ के लिए यथार्थ सर्वसमिका
        \[
        \norm{Hf_A}_2^2 = 4\ln A - 8 + \frac{8}{\sqrt A},
        \qquad\text{अतः}\qquad
        \frac{\norm{Hf_A}_2^2}{\norm{f_A}_2^2}
        = 4 - \frac{8 - 8A^{-1/2}}{\ln A} .
        \]
        सिद्ध कीजिए। $A = \eu^{10}$ पर मान निकालिए (अनुपात $\approx 1.790$) और
        टिप्पणी कीजिए: उच्चतम $2$ तक केवल $1/\ln A$ की गति से
        पहुँचा जाता है --- अर्थात् कोई इष्टतम नियतांक संख्यात्मक
        दृष्टि से लगभग अदृश्य हो सकता है।
\end{enumerate}
\end{problem}
